\chapter{Iterazione 3: API Gateway, Web UI e NLU Engine}
\label{chap:iterazione3}

\section{Obiettivi e Integrazione di Sistema}

L'Iterazione 3 segna il passaggio da una collezione di microservizi backend disgiunti a una piattaforma applicativa integrata. L'obiettivo primario è stato chiudere il cerchio architetturale, fornendo all'utente un punto di accesso unico e un'interfaccia naturale per interagire con i dati finanziari.

Le attività si sono concentrate su tre pilastri fondamentali:
\begin{itemize}
    \item \textbf{Centralizzazione dell'accesso}: Introduzione di un API Gateway per astrarre la topologia dei microservizi e risolvere le problematiche di CORS.
    \item \textbf{Interfaccia Utente (Web UI)}: Sviluppo di un frontend server-side (Thymeleaf) integrato con le API di riconoscimento vocale del browser.
    \item \textbf{Intelligenza Semantica (NLU)}: Implementazione di un motore deterministico per la traduzione del linguaggio naturale in comandi transazionali strutturati.
\end{itemize}

\section{API Gateway: Centralizzazione del Routing}

Nelle iterazioni precedenti, i servizi rispondevano su porte diverse (\texttt{8081}, \texttt{8082}), costringendo il client a conoscere la topologia della rete. Per disaccoppiare il frontend dal backend, è stato introdotto il microservizio \texttt{jarfin-api-gateway}, basato su Spring Cloud Gateway.

\subsection{Configurazione Dichiarativa delle Rotte}
Il Gateway espone un unico entry-point sulla porta 8080 e instrada le richieste ai servizi interni tramite regole definite nel file \texttt{application.yml}.

\begin{lstlisting}[language=yaml, label={lst:gateway_yml}]
server:
  port: 8080
spring:
  cloud:
    gateway:
      routes:
        # Rotta verso il Core Accounting
        - id: accounting-service
          uri: http://localhost:8081
          predicates:
            - Path=/api/transactions/**

        # Rotta verso la Business Intelligence
        - id: analytics-service
          uri: http://localhost:8082
          predicates:
            - Path=/api/analytics/**

        # Rotta verso il Frontend (Web UI)
        - id: web-ui
          uri: http://localhost:8083
          predicates:
            - Path=/**
\end{lstlisting}

Questa configurazione implementa il pattern \textbf{Reverse Proxy}: il browser dialoga esclusivamente con \texttt{localhost:8080}, ignorando l'esistenza dei servizi sottostanti. Questo semplifica la gestione della sicurezza e centralizza il logging delle richieste.

\section{Web UI e Riconoscimento Vocale Ibrido}

L'interfaccia utente è stata sviluppata utilizzando \textbf{Thymeleaf} e \textbf{Bootstrap 5}, servita dal microservizio \texttt{jarfin-web-ui}.
Una scelta architetturale critica ha riguardato l'implementazione del riconoscimento vocale (Speech-to-Text).

\subsection{Strategia "Edge AI" (Client-Side Processing)}
Invece di inviare flussi audio pesanti al server (che richiederebbero larghezza di banda e costose GPU per la trascrizione), si è optato per l'utilizzo delle \textbf{Web Speech API} native del browser.
La trascrizione avviene localmente sul dispositivo dell'utente (Edge), e solo il testo risultante viene inviato al backend.

Nel file \texttt{index.html}, la logica JavaScript gestisce il ciclo di vita del microfono e l'attivazione tramite "Wake Word" ("Johnny"):

\begin{lstlisting}[language=JavaScript, label={lst:js_stt}]
recognition.onresult = (event) => {
    // Concatenazione risultati parziali
    let transcript = "";
    for (let i = event.resultIndex; i < event.results.length; ++i) {
        if (event.results[i].isFinal) 
            transcript += event.results[i][0].transcript.toLowerCase();
    }
    
    // Rilevamento Wake Word locale
    if (!isJarfinActivated && WAKE_WORDS.some(w => transcript.includes(w))) {
        activateJohnny(); // Attiva la UI di ascolto
    } 
    // Invio comando al backend
    else if (isJarfinActivated) {
        sendToBackend(transcript.trim());
    }
};
\end{lstlisting}

Questa architettura ibrida garantisce \textbf{latenza zero} nell'attivazione e maggiore \textbf{privacy}, poiché l'audio raw non lascia mai il dispositivo dell'utente.

\newpage

% SONO ARRIVATO QUA

\section{Natural Language Understanding (NLU) Engine}

Il cuore intelligente del sistema è il servizio \texttt{NaturalLanguageService}. A differenza degli approcci basati su LLM (Large Language Models), che sono probabilisticamente corretti ma lenti e costosi, JarFin adotta un motore **deterministico basato su regole**.

Questa scelta garantisce:
\begin{itemize}
    \item \textbf{Velocità}: Elaborazione in $O(N)$.
    \item \textbf{Affidabilità}: A parità di input, l'output è sempre identico.
    \item \textbf{Interpretabilità}: Ogni estrazione è giustificata da una regola regex o keyword.
\end{itemize}

\subsection{Implementazione: Regex e Keyword Mapping}
Il servizio utilizza espressioni regolari pre-compilate per l'estrazione di importi e ID, e una mappa hash per la classificazione semantica.

\begin{lstlisting}[language=Java, caption={Costanti e Pattern nel NaturalLanguageService}, label={lst:nlu_constants}]
private static final Pattern NUMBER_PATTERN = 
    Pattern.compile("(\\d+([.,]\\d{1,2})?)"); // Supporta 10.50 e 10,50
    
private static final Pattern DELETE_CMD = 
    Pattern.compile("(elimina|cancella|rimuovi|togli)", Pattern.CASE_INSENSITIVE);

private static final Map<String, String> CATEGORY_MAP = new HashMap<>();
static {
    CATEGORY_MAP.put("mcdonald", "Ristorante");
    CATEGORY_MAP.put("benzina", "Trasporti");
    // ... mappatura estesa di oltre 50 keyword
}
\end{lstlisting}

\section{Design in Piccolo: Analisi dell'Algoritmo NLU}

L'algoritmo di parsing (\texttt{parse}) trasforma una stringa naturale non strutturata in un oggetto \texttt{ParsedTransaction} pronto per il database.

\subsection{Pseudocodice Formale}
L'algoritmo segue una pipeline sequenziale di estrazione e pulizia.

\begin{lstlisting}[mathescape=true, caption={Pseudocodice Parser NLU}, label={lst:pseudocode_nlu}, basicstyle=\small\ttfamily]
Algoritmo ParseNaturalLanguage(text):
    Input: Stringa $S$ (comando utente)
    Output: Oggetto $T$ (Transazione strutturata) o NULL

    // 1. Pre-validazione
    IF $S$ is NULL OR length($S$) < 2 OR contains(STOP_WORDS, $S$) THEN
        RETURN NULL
    END IF
    $S_{lower} \leftarrow$ ToLowerCase($S$)

    // 2. Identificazione Intento (Command Classification)
    IF Match($P_{delete}$, $S$) THEN 
        $T.command \leftarrow$ DELETE
        $T.targetId \leftarrow$ ExtractId($S$)
    ELSE IF Match($P_{update}$, $S$) THEN
        $T.command \leftarrow$ UPDATE
        $T.targetId \leftarrow$ ExtractId($S$)
    ELSE
        $T.command \leftarrow$ CREATE
        $T.date \leftarrow$ CurrentDate()
    END IF

    // 3. Estrazione Importo (Regex)
    Match $M_{num} \leftarrow$ Find($P_{number}$, $S$ without IDs)
    IF $M_{num}$ found THEN
        $T.amount \leftarrow$ ParseDecimal($M_{num}.value$)
    END IF

    // 4. Classificazione Categoria (Keyword Matching)
    FOR each ($Key$, $Cat$) in $Map_{categories}$ DO:
        IF $S_{lower}$ contains $Key$ THEN
            $T.category \leftarrow Cat$
            BREAK // Match trovato (Priority First)
        END IF
    END FOR

    // 5. Pulizia Descrizione (Stop-word Removal)
    $T.description \leftarrow$ CleanText($S_{lower}$) 
    IF $T.description$ is empty THEN
        $T.description \leftarrow T.category$ OR "Generale"
    END IF

    RETURN $T$
\end{lstlisting}

\subsection{Analisi della Complessità Computazionale}
Analizziamo il costo temporale $T(n)$ dell'algoritmo, dove $n$ è la lunghezza della stringa di input.

\begin{enumerate}
    \item \textbf{Regex Matching}: L'applicazione di pattern regex semplici (senza backtracking annidato o ricorsivo) su una stringa di lunghezza $n$ ha complessità lineare $O(n)$.
    \item \textbf{Keyword Matching}: L'algoritmo itera su una mappa di $K$ parole chiave. Per ogni parola chiave, l'operazione `contains` ha un costo proporzionale alla lunghezza della stringa.
    \[ T_{match} \approx K \cdot O(n) \]
    Poiché $K$ è una costante di configurazione (circa 50 entry) e non dipende dall'input utente, asintoticamente questo termine si riduce a $O(n)$.
    \item \textbf{Clean Text}: La rimozione delle stop-word avviene tramite una serie di \texttt{replaceAll}, ciascuno dei quali scorre la stringa una volta.
    \[ T_{clean} = \sum_{regex} O(n) \approx C \cdot O(n) \]
\end{enumerate}

\textbf{Complessità Totale}:
\[ T(n) = O(n) + K \cdot O(n) + C \cdot O(n) \]
\[ T(n) = O(n) \]

L'algoritmo è **Lineare**. Questo garantisce che il tempo di elaborazione sia trascurabile (< 1ms) per frasi di lunghezza standard (es. 10-100 caratteri), rendendo l'esperienza utente percepita come "istantanea".

\section{Orchestrazione: CommandController}

Il \texttt{CommandController} funge da orchestratore backend. Riceve il testo dal frontend, invoca il servizio NLU e instrada il comando strutturato verso l'API Gateway.

Il metodo \texttt{processVoiceCommand} gestisce la logica di fallback e l'invio al Core Accounting:

\begin{lstlisting}[language=Java, caption={Processamento comando vocale}, label={lst:command_controller}]
@PostMapping("/process")
public ResponseEntity<?> processVoiceCommand(@RequestBody Map<String, String> payload) {
    String text = payload.get("text");
    ParsedTransaction parsed = nluService.parse(text);

    // Instradamento dinamico basato sul tipo di comando
    switch (parsed.getCommandType()) {
        case CREATE:
            // Fallback: se la categoria non e' dedotta, usa "Altro"
            if (parsed.getCategory() == null) parsed.setCategory("Altro");
            if (parsed.getAmount() == null) parsed.setAmount(BigDecimal.ZERO);
            
            // Chiamata all'Accounting Service via Gateway
            restTemplate.postForEntity(gatewayUrl + "/api/transactions", 
                                       parsed, Object.class);
            break;
            
        case DELETE:
            // Logica di eliminazione transazione
            break;
    }
    return ResponseEntity.ok(Map.of("message", "Salvato"));
}
\end{lstlisting}

\section{Strategia di Validazione e Testing}

Allo stato attuale dello sviluppo, la validazione dell'Iterazione 3 è in fase di completamento. Di seguito viene illustrata la **Strategia di Test** progettata per garantire la robustezza dei nuovi componenti integrati.

\subsection{Unit Testing: NLU Engine}
L'obiettivo primario è verificare la precisione dell'algoritmo di parsing. Sono previsti test parametrizzati con JUnit 5 per coprire i seguenti scenari:
\begin{itemize}
    \item \textbf{Scenario Base}: Input "Ho speso 10 euro per pizza" $\rightarrow$ \texttt{Amount=10.00}, \texttt{Category=Ristorante}.
    \item \textbf{Edge Case}: Input con separatori decimali diversi ("10.50", "10,50") o valute miste.
    \item \textbf{Comandi Complessi}: "Elimina la transazione 42" $\rightarrow$ \texttt{Command=DELETE}, \texttt{TargetId=42}.
    \item \textbf{Input "Sporchi"}: Frasi lunghe con molte stop-word per verificare l'efficacia del metodo \texttt{cleanText}.
\end{itemize}

\subsection{Integration Testing: Gateway e Controller}
Per validare l'orchestrazione, verranno eseguiti test di integrazione che coinvolgono il container del Gateway:
\begin{itemize}
    \item \textbf{Routing Test}: Verifica che le chiamate a \texttt{localhost:8080/api/transactions} vengano correttamente inoltrate al servizio Accounting.
    \item \textbf{Flow Test}: Simulazione di una chiamata POST al \texttt{CommandController} con un payload testuale, verificando che l'Accounting Service riceva correttamente l'oggetto JSON strutturato.
\end{itemize}

Questa fase di Quality Assurance, attualmente in corso, certificherà la stabilità della piattaforma prima del rilascio finale.

\section{Conclusione Iterazione 3}

L'Iterazione 3 ha completato l'architettura JarFin.
L'API Gateway ha fornito un layer di astrazione essenziale per la sicurezza e il routing.
Il motore NLU, grazie alla sua natura algoritmica $O(n)$, offre un'alternativa robusta e performante all'IA generativa per compiti strutturati, garantendo un'esperienza utente fluida e priva di latenze significative.
